\documentclass[11pt]{article}
\usepackage[margin=0.78in]{geometry}
\usepackage{amsmath,amssymb,bm}
\usepackage{booktabs}
\usepackage{enumitem}
\usepackage{hyperref}
\hypersetup{colorlinks=true, linkcolor=black, urlcolor=black}
\setlength{\parindent}{0pt}
\setlength{\parskip}{0.55em}

\title{Cyrus Guided Notes: MIT Underactuated Robotics}
\author{MIT course pack}
\date{2026-05-15}

\begin{document}
\maketitle

\section*{Course source}
\textbf{Official source:} \url{https://underactuated.mit.edu/}

\textbf{Why this course is in the system.} Entry point for robot dynamics, optimal control, trajectories, stability, and nonlinear systems.

\textbf{Learning output.} Build inverted-pendulum, phase-portrait, and trajectory-optimization interactions on the web.

\section*{Prerequisite checkpoint}
\begin{itemize}[leftmargin=1.4em]
  \item Dynamics
  \item Optimization
  \item State-space control
\end{itemize}

\section*{Formula checkpoint}
Do not memorize these first. Read each formula as a sentence: what changes, what is observed, what is optimized, and what output it produces.
\begin{align*}
  M(\bm{q})\ddot{\bm{q}}+C(\bm{q},\dot{\bm{q}})\dot{\bm{q}}+\bm{g}(\bm{q})=B\bm{u}\\
  J=\sum_{k=0}^{N}\ell(\bm{x}_k,\bm{u}_k)\\
  V(\bm{x})>0,\quad \dot{V}(\bm{x})<0\\
  \bm{x}_{k+1}=f(\bm{x}_k,\bm{u}_k)
\end{align*}

\section*{Visual page}
Use an inverted pendulum phase portrait to show underactuation and stabilization.

\section*{GoodNotes pages}
\begin{itemize}[leftmargin=1.4em]
  \item Page MIT-R001: underactuated dynamics
  \item Page MIT-R002: trajectory optimization
\end{itemize}

\section*{Web interaction}
A pendulum slider for angle and control effort.

\section*{Obsidian and Notion}
\begin{tabular}{p{0.22\linewidth}p{0.68\linewidth}}
\toprule
Layer & Output \\
\midrule
Obsidian & MIT Robotics -> underactuated robotics \\
Notion & Robot-control concepts and experiment state. \\
\bottomrule
\end{tabular}

\section*{First study move}
Open the official source, copy the prerequisite checkpoint into GoodNotes, then write one formula in your own words before watching or reading more.

\end{document}
